% from ega0-16-fr.tex
\oldpage[0\textsubscript{IV-1}]{37}

\begin{proposition}[16.5.5]
\label{0.16.5.5-fr}
Soient $A$ un anneau local noethérien, $\mathfrak{m}$ son idéal maximal, $M$
un $A$-module de type fini $\neq0$, $x$ un élément de $\mathfrak{m}$.
Supposons que $M$ soit un module de Cohen-Macaulay~; alors les conditions
suivantes sont équivalentes~:
\begin{enumerate}
  \item[a)] $x$ est $M$-régulier~;
  \item[b)] $\dim(M/xM)=\dim(M)-1$~;
  \item[c)] $x$ n'appartient à aucun élément minimal de
  $\operatorname{Supp}(M)$.
\end{enumerate}
En outre, $M/xM$ est alors un module de Cohen-Macaulay.
\end{proposition}

\begin{proof}
On a vu (16.5.4) que tous les éléments $\mathfrak{p}$ de
$\operatorname{Ass}(M)$ sont minimaux dans $\operatorname{Supp}(M)$ et que
$\dim(A/\mathfrak{p})=\dim(M)$~; l'équivalence de a) et c) résulte donc de
Bourbaki, \emph{Alg.~comm.}, chap.~IV, \S~1, no~1, cor.~2 de la prop.~2~;
l'équivalence de b) et c) résulte de (16.3.4). Le fait que $M/xM$ soit un
module de Cohen-Macaulay résulte enfin de (16.4.10, (i)).
\end{proof}

\begin{corollary}[16.5.6]
\label{0.16.5.6-fr}
Sous les hypothèses de (16.5.5), soit $(x_i)_{1\leq i\leq r}$ une suite
d'éléments de $\mathfrak{m}$. Les conditions suivantes sont équivalentes~:
\begin{enumerate}
  \item[a)] La suite $(x_i)$ est $M$-régulière.
  \item[b)] On a
  $\dim(M/(x_1M+\cdots+x_rM))=\dim(M)-r$.
  \item[c)] Les $x_i$ font partie d'un système de paramètres de $M$.
\end{enumerate}
En outre, lorsque ces conditions sont satisfaites,
$N=M/(x_1M+\cdots+x_rM)$ est un module de Cohen-Macaulay, donc, pour tout
idéal $\mathfrak{p}\in\operatorname{Ass}(N)$, on a
$\dim(A/\mathfrak{p})=\dim(M)-r$.
\end{corollary}

\begin{proof}
On sait déjà que b) et c) sont équivalentes (16.3.6) et que a) entraîne c)
(16.4.1) sans hypothèse sur $M$. Pour voir que c) implique a) et prouver la
dernière assertion de l'énoncé, raisonnons par récurrence sur $r$, l'assertion
étant triviale pour $r=0$. Puisque $x_1$ fait partie d'un système de
paramètres, $\dim(M/x_1M)=\dim(M)-1$ (16.3.6), donc $x_1$ est $M$-régulier et
$M/x_1M$ est un module de Cohen-Macaulay par (16.5.5). Comme la suite
$(x_i)_{2\leq i\leq r}$ fait partie d'un système de paramètres pour
$P=M/x_1M$, l'hypothèse de récurrence montre que $(x_i)_{2\leq i\leq r}$ est
une suite $P$-régulière et que
$P/(x_2P+\cdots+x_rP)=N$ est un module de Cohen-Macaulay~; en outre
$(x_i)_{1\leq i\leq r}$ est $M$-régulière.
\end{proof}

\begin{remark}[16.5.7]
\label{0.16.5.7-fr}
Il résulte de (16.5.6) qu'il y a identité entre les \emph{systèmes de
paramètres pour $M$} et les \emph{suites $M$-régulières maximales d'éléments de
$\mathfrak{m}$} lorsque $M$ est un module de Cohen-Macaulay. Inversement, il
résulte de (16.4.1) que si un $A$-module de type fini non nul est tel qu'il y
ait une suite $M$-régulière qui soit un système de paramètres pour $M$, alors
$M$ est un module de Cohen-Macaulay. La dernière propriété énoncée dans
(16.5.6) \emph{caractérise aussi} les modules de Cohen-Macaulay~:
\end{remark}

\begin{proposition}[16.5.8]
\label{0.16.5.8-fr}
Soient $A$ un anneau local noethérien, $M$ un $A$-module de type fini. On
suppose que pour toute suite $(x_i)_{1\leq i\leq r}$ d'éléments de $A$
faisant partie d'un système de paramètres pour $M$, et tout idéal premier
$\mathfrak{p}$ associé à
$N=M/(x_1M+\cdots+x_rM)$, on ait
$\dim(A/\mathfrak{p})=\dim(N)$. Alors $M$ est un module de Cohen-Macaulay.
\end{proposition}

\begin{proof}
Soit $(y_i)_{1\leq i\leq n}$ un système de paramètres pour $M$~; il suffit de
montrer que la suite $(y_i)$ est $M$-régulière. Raisonnons par récurrence sur
$n=\dim(M)$, la proposition étant triviale pour $n=0$. Par hypothèse
(appliquée avec $r=0$) les idéaux premiers $\mathfrak{p}_j$ associés à $M$
sont tous tels que $\dim(A/\mathfrak{p}_j)=\dim(M)$~; comme $y_1$ fait partie
d'un système de paramètres
